\section{Introduction}
A fixed-wing aircraft must continue moving while its planner chooses the next maneuver. A trajectory that is collision-free over a short horizon can terminate with insufficient room for the subsequent turn or climb. Motion-primitive planners provide a structured way to handle vehicle dynamics, but selecting a primary trajectory solely by immediate task cost can overlook this loss of maneuvering space.

Two-phase fixed-wing planning combines global route guidance with local maneuver selection~\cite{tan2022twophase,tan2026twophase}. Within this hierarchy, a terminal continuation test evaluates a second control segment from the complete endpoint state of each primary. The additional look-ahead introduces a computational problem: checking $M$ successors for each of $N$ primaries requires up to $NM$ dynamic rollouts. A limited validation budget must therefore be allocated among candidates that differ in both their likelihood of admitting a continuation and the effort required to establish that result.

Task-cost ordering tends to concentrate checks on similar terminal motions. Maneuver diversity spreads the search across turning directions and altitudes, but geometry can make some of these alternatives unproductive. An endpoint near an open turning region and one facing a wall may be equally distant from the previously selected maneuvers. This motivates an environment-conditioned schedule that learns which terminal states are promising while preserving exploration across distinct maneuvers.

We develop a learning-guided continuation search for a fixed-wing primitive planner. A compact neural network receives the initial and terminal dynamic states, target information and a terminal-frame clearance representation. It predicts whether a finite continuation witness exists and the work required by the incremental validator. The predictions prioritize explicit checks, and a hybrid schedule interleaves learned ranking with physically normalized maneuver coverage. The primary generator, task objective and trajectory acceptance tests are shared with the baselines.

The contributions are threefold. First, we formulate learned continuation-query selection under candidate-count and integration-step budgets, using labels obtained directly from a fixed-wing dynamic validator. Second, we combine a geometry-conditioned neural ranking with maneuver-space exploration and incremental rejection of infeasible continuations. Third, we evaluate the learned component on independent map splits through geometry ablation, matched-state budget comparisons and closed-loop execution with nominal dynamics, model mismatch, wind and command delay. An additional urban Gazebo flight evaluates online execution with onboard three-dimensional lidar acquisition. The analysis separates candidate prediction accuracy from validation yield and mission completion. The study isolates the local planning layer: every method receives the same target and obstacle map.

\section{Related Work}
\subsection{Fixed-Wing Planning and Terminal Maneuvers}
Motion primitives and hierarchical route planning provide a computationally structured approach to fixed-wing navigation~\cite{tan2022twophase,tan2026twophase}. The implementation in this paper builds on the public three-dimensional local primitive planner in~\cite{hkpolyu2026code}, retaining its point-mass dynamics and actuator responses.

Terminal constraints in model predictive control connect finite-horizon optimization with closed-loop feasibility~\cite{mayne2000mpc}. Inevitable collision states characterize states from which all future collision-free motions have been lost~\cite{fraichard2004ics}. For fixed-wing navigation, Lim et al.~\cite{lim2024safe} connect trajectories to collision-free periodic circles, whose repeatability supports their safe-state construction. FASTER~\cite{tordesillas2022faster} maintains a backup trajectory within known free space while allowing a task trajectory to explore farther.

Our continuation query asks whether a primary endpoint admits a feasible successor from a specified control library over a specified duration. This finite witness fits the existing fixed-wing generator and provides a measurable validation task. The scheduling method allocates these queries; the periodic terminal maneuver and invariant-set constructions above address a different horizon of feasibility.

% Replace the former Backup Trajectories and Learned Candidate Libraries subsection
% with this text, and retain the separate terminal-set and diversity discussions.
\subsection{Learned Primitives and Validation Effort}
YOPO~\cite{lu2024yopo} learns trajectory parameters and scores for a spatially anchored candidate set, while YOPOv2-Tracker~\cite{lu2026yopov2} extends this representation to integrated tracking and navigation. ORACLE~\cite{nguyen2022oracle} predicts collision costs for predefined velocity--steering action sequences from depth images and estimated motion, and combines these predictions with goal information to select a primitive. DPNet~\cite{zhang2026dpnet} predicts structural dead ends and prunes a compact trajectory library. These approaches establish the role of learned environment-dependent scores in selecting among multiple local maneuvers.

Learning can also target the computation performed by a planner. SaIL~\cite{bhardwaj2017sail} trains a node-expansion policy by imitating a fully informed search oracle, with search effort as its objective. Bayesian active edge evaluation~\cite{choudhury2018active} uses correlations among edge-validity outcomes to choose informative checks in a path library. Fastron~\cite{das2017fastron} instead learns a configuration-space collision proxy. These methods address different parts of the planning pipeline: search expansion, evaluation order, and the collision test itself.

We learn the priority of terminal continuation queries for a dynamically propagated fixed-wing primitive library. Each query asks whether the terminal state admits a feasible successor and how many integration steps the incremental validator is likely to require. The two quantities provide joint supervision from feasibility and computational-effort labels. The influence of predicted work on ranking is selected using validation maps. The online planner combines this learned priority with maneuver-space coverage, and accepts a candidate after evaluating the corresponding continuation under the dynamics and geometric constraints. The comparison therefore concerns how much useful continuation evidence is obtained within a prescribed budget.


\subsection{Maneuver Coverage}
Farthest-first selection spreads representatives across a metric space~\cite{gonzalez1985clustering}. We use this principle to diversify terminal headings and altitudes, normalized by nominal aircraft maneuver capability. Such physical scaling prevents small altitude variations in a narrow flight band from receiving the same importance as a substantial heading change. The hybrid schedule couples this coverage with learned geometric information, and the experiments isolate the contribution of each component.
